\section{Conclusion}\label{sec:conclusion}

This paper presented an agentic fraud detection system built on the MAPE-K architecture. The system integrates four key components: data collection and LLM annotation from social media, CTGAN-based data augmentation to handle class imbalance, adversarial training for robustness, and autonomous adaptation using MAPE-K agents.

Our experiments showed strong results. The system achieves F1 score of 0.9911 and ROC-AUC of 0.9966 on real fraud data. The robustness against adversarial attacks is 98.34 percent, demonstrating that the model can handle attempts to evade detection. The MAPE-K architecture successfully detects concept drift and triggers automatic retraining, achieving F1 of 0.92 even under changing fraud patterns.

The key contributions of this work are a complete agentic fraud detection pipeline that handles data collection, augmentation, adversarial training, and autonomous adaptation, a practical implementation of MAPE-K for fraud detection that demonstrates measurable learning adaptation, evidence that combining CTGAN with adversarial training produces highly robust models, and strong experimental results validated on real-world data.

For practitioners, this system offers a deployable solution that reduces manual oversight through automatic adaptation. For researchers, the integration of multiple techniques into a coherent framework provides a foundation for further work.

Future directions include exploring more advanced model architectures, testing on additional datasets from different sources, and implementing more sophisticated adversarial training methods. The MAPE-K framework could also be extended with better planning algorithms to optimize retraining decisions.

In conclusion, our agentic fraud detection system demonstrates how intelligent agents can work together to handle the complex challenges of modern fraud detection. By combining established techniques with autonomous adaptation, we create a robust and practical solution for real-world fraud detection needs.